% Source: https://artofproblemsolving.com/community/c3h1383698_midpoints_and_angle_bisector_problem
% Stable identity: aops-topic-1383698-q01
By the midsegment theorem, $MN\parallel BC$ and $MN=\dfrac{1}{2}BC=30$. Since $MN\parallel BC$, the alternate angles $\angle MPB$ and $\angle PBC$ are equal, and $\angle PBC=\angle MBP$ because $BT$ bisects $\angle ABC$. So $\angle MPB=\angle MBP$ and triangle $MBP$ is isosceles:
\[
MP=MB=\dfrac{45}{2}=22.5 .
\]

\begin{asy}
usepackage("amsmath");
size(6.5cm);
defaultpen(fontsize(10pt));
pair ft(pair P,pair U,pair V){return U+dot(P-U,V-U)/dot(V-U,V-U)*(V-U);}
path ra(pair P,pair u,pair v,real s=0.2){return P+s*unit(u)--P+s*unit(u)+s*unit(v)--P+s*unit(v);}
pair B=(0,0),C=(60,0);
real ax=(45^2-70^2+60^2)/120; pair A=(ax,sqrt(45^2-ax^2));
pair M=(A+B)/2,Nn=(A+C)/2;
pair T=A+(45/105)*(C-A);
pair P=extension(B,T,M,Nn);
draw(A--B--C--cycle,linewidth(0.9));
draw(M--Nn,linewidth(0.8)); draw(B--T,blue);
draw(arc(B,7,0,degrees(T-B)),blue); draw(arc(B,8.2,degrees(T-B),degrees(A-B)),blue);
dot(M);dot(Nn);dot(P);dot(T);
label("$A$",A,N); label("$B$",B,SW); label("$C$",C,SE); label("$M$",M,W); label("$N$",Nn,E);
label("$T$",T,NE); label("$P$",P,dir(255));
label("$60$",(B+C)/2,S);
label("$22.5$",(M+B)/2,NW,fontsize(8pt));
label("$22.5$",(M+P)/2,N,fontsize(8pt)); label("$7.5$",(P+Nn)/2,S,fontsize(8pt));
label("\mbox{}",(min(currentpicture,user=true)+max(currentpicture,user=true))/2,black); // keeps the text after the figure black when LaTeX typesets the labels inline
\end{asy}

(The figure is drawn with $AC=70$; the answer does not depend on $AC$.) Then $PN=30-22.5=7.5$, and
\[
\dfrac{MP}{PN}=\dfrac{22.5}{7.5}=\boxed{3}.
\]
